\section{Discussion}\label{sec:discussion}

\subsection{What the evidence supports}\label{sec:disc-support}

Under a protocol that keeps uniform data out of training, apart from Coat's shipped propensities, and selects on validation data, the ranking of debiasing methods depends on evaluation choices in ways that matter at the scale of the differences being compared. Selecting on test data inflates the Coat scores by up to $0.030$ nDCG@5 and reverses the order of iALS and LightGCN. The differences among the leading Coat methods are of the same order (Table~\ref{tab:T4}), so a ranking of those methods made on test data is not reliable without the validation step.

Graph-specific debiasing adds little over an equally tuned backbone on Coat, Yahoo! R3 and KuaiRand. The largest gain over LightGCN on these three datasets is $+4.10$ ($\times10^{-3}$) nDCG, for r-AdjNorm on Yahoo! R3, and the Nadeau-Bengio correction removes its significance. KuaiRec is different. NAVIP and DR-LightGCN add $+10.41$ and $+13.17$ ($\times10^{-3}$) there, although r-AdjNorm adds only $+0.22$, and a change of training loss without any debiasing adds $+19.59$. A reader comparing debiasers on KuaiRec therefore needs the pointwise backbone as the reference, not only the BPR one.

On Coat and KuaiRec the additive prior is the ingredient with the largest effect that we tested. Added as a score offset to LightGCN, it gives $+31.98$ [20.65] on Coat and $+21.22$ [2.86] on KuaiRec, and it does not change Yahoo! R3 or KuaiRand. The evidence does not show that the inverse-propensity weighting matters: on KuaiRec the unweighted prior is significantly better than the weighted one, and on Coat it is numerically higher without a significant difference. A random permutation of the item term removes the KuaiRec gain, but on Coat the same permutation does not significantly change the score, so the item term is established on KuaiRec and not on Coat. These results do not establish why the prior helps. The reading that the item term uses exposed negatives that a BPR objective discards is a hypothesis, and the prior-alone comparisons do not test it.

On Yahoo! R3 and KuaiRand the weights on the sampled negatives carry the gain. On Yahoo! R3 a shuffled weight vector removes it; no shuffled control was run on KuaiRand. Popularity weights are significantly below LightGCN on Yahoo! R3, so the weights are not a popularity reweighting. On KuaiRand, however, the popularity prior alone outperforms the IPS-fitted prior alone by $+56.41$ [5.19], so the simple non-personalised baseline is a serious competitor there, and a study that omits it would overstate the debiased prior.

On KuaiRec, training-free doubly robust propagation has the highest mean in Table~\ref{tab:T4}, and no trained method or refinement in the study exceeds it. The walk correction that makes its propagation exactly unbiased on synthetic logs does not change accuracy on the real logs at matched configurations. For ranking, then, the bias that the correction removes is not the limiting factor in these data.

Finally, the backbone matters as much as the debiaser. A wider iALS grid on Coat raises the stored iALS by $+12.92$ [8.01] (p$_\mathrm{Holm}$ = 0.003), against $+16.82$ [18.04] for the offset model over the same stored iALS. An evaluation that tunes the proposed method more than its baselines can produce a gap of that size from budget alone.

\subsection{Implications for evaluating debiasing methods}\label{sec:disc-implications}

The results support the following practices. Each item rests on a result above.

\begin{enumerate}
\item \textbf{Keep uniform data out of training, and state where every nuisance comes from.} Coat's shipped propensities came from the uniform sample. Refitting them with the estimator used on the other datasets changes the Coat prior results by at most $+2.63$ [5.27] (Table~\ref{tab:T10}). The check is cheap, and it should be reported whenever the propensities are not fitted on the training log.
\item \textbf{Select per split on validation data, and report the test-oracle gap.} The gap is $+0.0097$ to $+0.0302$ on Coat, and it reverses the order of iALS and LightGCN there (Section~\ref{sec:res-eval}).
\item \textbf{Report the budget, and compare at matched budget.} Grid sizes in these experiments range from 6 to 324 configurations. The sign of the DRUP comparison with DR-JL on KuaiRec changes in 78.4\% of six-configuration subgrids (Table~\ref{tab:T7}).
\item \textbf{Include the simple baselines as mandatory rows:} the IPS imputation alone, a tuned iALS with its wider grid, pointwise LightGCN, the popularity ranker, and the backbone with the debiased prior. Several of these changed a conclusion here: pointwise LightGCN on KuaiRec ($+19.59$ [4.55] over the BPR backbone), popularity against the prior alone on KuaiRand, and the wider iALS grid on Coat.
\item \textbf{Retrain the selected configurations with several seeds.} On Coat, the difference between \qp{} and iALS is $+20.02$ [19.74] (p$_\mathrm{Holm}$ = 0.0469) in the main run and $+3.39$ [20.70] (p$_\mathrm{Holm}$ = 0.747) over seeds 1 to 3.
\item \textbf{Report the minimum detectable difference with every claim of no difference, and test equivalence with a margin fixed in advance.} On Coat the minimum detectable difference for the primary comparisons is $0.012$ to $0.043$, which is large relative to the method differences in Table~\ref{tab:T4}.
\item \textbf{Use a correction for overlapping splits, and report it next to the standard test.} Under the Nadeau-Bengio correction, no graph-debiaser comparison (P1) is significant on any dataset.
\item \textbf{Publish the per-split validation and test results.} The selection analyses here, including the inflation, the subgrid stability and the bootstrap ranks, can be redone only from per-split values.
\end{enumerate}

\subsection{Limitations and threats to validity}\label{sec:disc-limits}

\begin{enumerate}
\item \textbf{Power on Coat.} Coat has 225 evaluated users. The minimum detectable differences of the primary comparisons range from $0.012$ to $0.043$. Coat results are therefore statements of the form ``no detectable difference, and differences larger than the stated margin are excluded'', never statements of equivalence.
\item \textbf{Provisional equivalence margins.} The TOST margins ($\pm0.010$ on Coat and Yahoo! R3, $\pm0.003$ on KuaiRand and KuaiRec) were declared before the analysis. They have not yet been checked against published effect sizes, and they should be read as provisional.
\item \textbf{Our implementations.} The graph baselines NAVIP, r-AdjNorm, SimGCL and MACR, and the refinements, are our implementations. AdvInfoNCE is a simplified re-implementation: its adversary budget is implicit, it has no KL-ball projection, and its grid differs across datasets. The official AdvInfoNCE code and the PPD post-hoc method were not run. The negative results for these methods are therefore about our implementations under this protocol.
\item \textbf{The refinement family is ours.} The seven refinements (Section~\ref{sec:res-ingredient}) were developed as candidate methods before this study, and they are reported in full with their configuration counts. The family is not exhaustive, and the results do not show that no other refinement helps.
\item \textbf{Seeds are incomplete.} Seeds 1 to 3 were not run for the Coat prior rows. Seed checks cover the leading methods only. DRGP was run on Coat and KuaiRec only, and its KuaiRec and Coat stage-1 results have no seeds.
\item \textbf{Budgets are unequal and only partly matched.} Matched-budget differences are split means without tests. The sensitivity analysis draws random subgrids of the stored configurations, not new searches.
\item \textbf{Grid edges.} On KuaiRec, the popularity and shuffled priors select the lowest offset weight of their grid in all five splits, so the reported shuffled-prior harm is an upper bound under that grid. On KuaiRec the unweighted prior selects the largest regularisation value of its grid (1000). The wider iALS grid is not exhausted.
\item \textbf{Incomplete ablations.} The user-only prior was run only as a prior alone and not in a learner. LightGCN with estimated propensities was not run. On KuaiRec, only the additive prior was run in the offset-only LightGCN setting.
\item \textbf{Replay gaps.} Replaying each stored selection from its bank reproduces 151 of 154 method means within 0.002. Three do not. On replay the KuaiRec DR-LightGCN mean is $0.0033$ higher than the stored $0.6212$, the Yahoo! R3 DR-LightGCN mean $0.0058$ higher, and the Yahoo! R3 MACR mean $0.0060$ higher. The stored values are reproduced only with an early-stopping setting that the bank does not record. The KuaiRec and Yahoo! R3 DR-LightGCN values appear in Table~\ref{tab:T4}, and the KuaiRec value also enters the P1 comparison of Table~\ref{tab:T5}.
\item \textbf{Family dependence.} Several p-values change with the Holm family. The Yahoo! R3 r-AdjNorm comparison has p$_\mathrm{Holm}$ = 0.029 in the primary family and 0.073 over all comparisons with LightGCN. The Coat offset comparison has 0.013 in the primary family and 0.046 in the family of the prior-anatomy table. We report the family with each value.
\item \textbf{Overlapping splits.} Per-user tests over overlapping splits treat the log and the fitted nuisances as fixed, and they can be anti-conservative. The Nadeau-Bengio correction removes significance from every P1 comparison.
\item \textbf{Exposure semantics.} On Coat and Yahoo! R3, exposure is user self-selection. The derivation of the weights assumes that exposure is independent of the outcome given the covariates of the propensity model.
\item \textbf{Test design.} We report one K per dataset in the main text (nDCG@5 on Coat and Yahoo! R3, nDCG@10 on KuaiRand, nDCG@20 on KuaiRec). The alternative metric is computed on the stored selections only. KuaiRec has a fully observed test matrix with about 3,300 candidates per user, which differs from the other three datasets, and we have no test of which feature causes the difference.
\item \textbf{Unexplained differences.} Two fits of the imputation prior on KuaiRec give $0.6194$ and $0.6216$ (Table~\ref{tab:T4} and Table~\ref{tab:T9}). No paired test compares the prior alone with DRUP, so the $0.0011$ gap to DRUP is a difference of means.
\item \textbf{Synthetic evidence for the walk correction.} The bias results come from a synthetic MNAR model and a semi-synthetic KuaiRec study. Their gains do not carry over to the real logs at matched configurations.
\item \textbf{The mechanism is untested.} The item-term reading in Section~\ref{sec:disc-support} is a hypothesis. The results show that an item-level term fitted once on the exposed log carries the effect on KuaiRec. They do not show why.
\end{enumerate}
